\pdfminorversion=7%
\documentclass[aspectratio=169,mathserif,notheorems]{beamer}%
%
\xdef\bookbaseDir{../../bookbase}%
\xdef\sharedDir{../../shared}%
\RequirePackage{\bookbaseDir/styles/slides}%
\RequirePackage{\sharedDir/styles/styles}%
\toggleToGerman%
%
\subtitle{40.~Logisches Schema:~Beziehungen höheren Grades}%
%
\begin{document}%
%
\startPresentation%
%
\section{Einleitung}%
%
\begin{frame}[t]%
\frametitle{Einleitung}%
\begin{itemize}%
%
\only<-4>{%
\item Hier reproduzieren wir eines unserer frühen \glsShort{ERD}s mit einer ternäre Beziehung.%
}%
%
\item<2-> Die Entitätstypen \emph{Professor}, \emph{Student}, und \emph{Module} stehen miteinander in Beziehung und diese Beziehung hat sogar Attribute.%
%
\item<3-> Der Professor unterrichtet in Modul in einem bestimmten Semester.%
%
\item<4-> Die Studentin schreibt sich in das unterrichtete Modul in dem Semester ein.%
%
\item<5-> Das wollen wir nun in ein logisches Modell transferrieren.%
\end{itemize}%
\locateGraphic{}{width=0.6\paperwidth}{graphics/erdStudentModuleProf3}{0.2}{0.45}%
\end{frame}%
%
\section{Beispiel}%
%
%
\begin{frame}[t]%
\frametitle{Logisches Modell: Grundlagen}%
\begin{itemize}%
%
\item Bauen wir nun ein logisches Modell, dass zu diesem Szenario passt.%
%
\item<2-> Erst machen wir mal die einfachen Dinge.%
%
\item<3-> Aus den Entitätstypen werden wieder Tabellen mit Ersatz-Primärschlüsseln.%
%
\item<5-> Wir bekommen die Tabellen \sqlil{module}, \sqlil{professor} und \sqlil{student}.%
%
\item<6-> Damit das Beispiel etwas anschaulicher wird, fügen wir noch das Attribut \sqlil{name} zu den Tabellen \sqlil{student} und \sqlil{professor} hinzu und \sqlil{title} zu \sqlil{module}.%
%
\end{itemize}%
\locateGraphic{-3}{width=0.6\paperwidth}{graphics/erdStudentModuleProf3}{0.2}{0.45}%
\locateGraphic{4-5}{width=0.9\paperwidth}{graphics/logicalStudentModuleProfc}{0.05}{0.6}%
\locateGraphic{6-}{width=0.9\paperwidth}{graphics/logicalStudentModuleProfb}{0.05}{0.6}%
\end{frame}%
%
%
\begin{frame}[t]%
\frametitle{Logisches Modell: Beziehungen}%
\begin{itemize}%
%
\item Nun wollen wir die ternäre Beziehung bauen.%
%
\item<2-> Dafür gibt es im Grunde zwei Möglichkeiten\only<-2>{.}\uncover<3->{:%
\begin{enumerate}%
\item Wir können eine einzelne neue Tabelle \sqlil{takes_place} bauen.\uncover<4->{ %
Diese bekommt drei Fremdschlüssel auf die Tabellen \sqlil{professor}, \sqlil{student} und \sqlil{module}, sowie eine Spalte~\sqlil{semester}.\uncover<5->{ %
Sie speichert dann alle Daten der Beziehung.%
}}%
%
\item<6-> Wir können eine Tabelle \sqlil{course} mit einem Ersatz-Primärschlüssel erstellen und benutzen sie, um die Instanzen von Modulen zu speichern.\uncover<7->{ %
Eine Modul-Instanz ist sozusagen, wenn ein Modul von einem Professor in einem Semester unterrichtet wird.\uncover<8->{ %
Die Tabelle hat daher Fremdschlüssel auf die Tabellen \sqlil{professor} und \sqlil{module} und die Spalte \sqlil{semester}.\uncover<9->{ %
Wir erstellen dann eine zweite Tabelle \sqlil{enrolls} mit zwei Fremdschlüsseln, eine auf die Tabelle \sqlil{student} und eine auf die Tabelle \sqlil{course}.\uncover<10->{ %
Sie speichert die \inQuotes{schreibt-sich-ein}~\inEN{enrolls}-Beziehung der Studenten zu den Modul-Instanzen.%
}}}}%
\end{enumerate}%
}%
%
\item<11-> Beide Lösungen sind machbar.%
\end{itemize}%
\end{frame}%
%
\begin{frame}%
\frametitle{Lösung}%
\begin{itemize}%
\item Beide Lösungen sind machbar.\uncover<2->{ Wir nehmen aber die zweite.}%
%
\item<3-> Die Grundidee ist, dass wir die Realisierung eines Moduls von einem Professor in einem Semester als Entität betrachten können.%
%
\item<4-> Diese Entität hat das Attribut \sqlil{semester}, könnte aber auch noch mehr Attribute haben.%
%
\item<5-> Die Studenten dann mit dieser \sqlil{course} Tabelle über eine weitere, separate Tabelle in Beziehung zu setzen ergibt Sinn.%
%
\item<6-> Dieser Ansatz hat zwei große Vorteile\only<-6>{.}\uncover<7->{:%
%
\begin{enumerate}%
%
\item<7-> Erstens vermeidet er Redundanz.\uncover<8->{ %
Nach der ersten Idee würden wir den Fakt dass \inQuotes{Frau Prof.\ Bibba} das Modul \inQuotes{Mathematics~101} unterrichtet mehrfach speichern.\uncover<9->{ %
Wir speichern diesen Fakt nun nur noch einmal, nämlich in der Tabelle \sqlil{course}.%
}}%
%
\item<10-> Mit der zweiten Methode können wir auch die Situation repräsentieren, das ein Professor das selbe Modul zweimal in einem Semester unterichtet.\uncover<11->{ %
Kann ja sein, dass das auf Grund von Zeit- oder Raumkonflikten notwendig ist.%
}%
\end{enumerate}%
}%
\end{itemize}%
\end{frame}%
%
\begin{frame}[t]%
\frametitle{Lösung}%
\begin{itemize}%
\item Wir zeichnen also ein neues \glsShort{ERD} mit dem \pgmodeler.%
%
\item<5-> Wir treffen folgende Annahmen\only<-5>{.}\uncover<6->{:%
%
\begin{enumerate}%
%
\only<-11>{%
\setcounter{enumi}{0}%
\item Jedes Modul kann beliebig oft als Kurs instantiiert werden.%
}%
%
\only<-12>{%
\setcounter{enumi}{1}%
\item<7-> Jeder Kurs gehört zu exakt einem Modul.%
}%
%
\only<-13>{%
\setcounter{enumi}{2}%
\item<8-> Jeder Professor kann beliebig viele Module unterrichten.%
}%
%
\only<-14>{%
\setcounter{enumi}{3}%
\item<9-> Jeder Kurs gehört zu exakt einem Professor.%
}%
%
\only<-14>{%
\setcounter{enumi}{4}%
\item<10-> Jede Studentin kann sich in beliebig viele Kurse einschreiben.%
}%
%
\only<-14>{%
\setcounter{enumi}{5}%
\item<11-> Jede \inQuotes{Einschreibung} ist exakt einer Studentin und exakt einem Kurs zugeordnet.%
}%
%
\only<-14>{%
\setcounter{enumi}{6}%
\item<12-> Beliebig viele Studenten können sich in einen Kurs einschreiben.%
}%
%
\setcounter{enumi}{7}%
\item<13-> Jede Studentin kann sich nicht mehr als einmal in einen Kurs einschreiben.\uncover<14->{ %
Wir machen die beiden Fremdschlüssel in \sqlil{enrolls} zu einem kombinierten Primärschlüssel und erzwingen so ihre \sqlil{UNIQUE}ness.}%
\end{enumerate}%
}%
\end{itemize}%
%
\locateGraphic{-1}{width=0.6\paperwidth}{graphics/erdStudentModuleProf3}{0.2}{0.3}%
\locateGraphic{2}{width=0.9\paperwidth}{graphics/logicalStudentModuleProfc}{0.05}{0.35}%
\locateGraphic{3}{width=0.9\paperwidth}{graphics/logicalStudentModuleProfb}{0.05}{0.35}%
\locateGraphic{4}{width=0.9\paperwidth}{graphics/logicalStudentModuleProf}{0.05}{0.3}%
\locateGraphic{5-}{width=0.5\paperwidth}{graphics/logicalStudentModuleProf}{0.25}{0.6}%
\end{frame}%
%
%
\section{Generiertes SQL}%
%
\begin{frame}[t]%
\frametitle{Datenbank erstellen}%
\begin{itemize}%
\item Erstellen wir nun die Datenbank und Tabellen mit den Skripten, die der \pgmodeler\ für uns generiert hat.%
%
\item<2-> Fangen wir mit der Datenbank an.
\end{itemize}%
%
%
\gitLoadAndExecSQL{teaching_1:01_teaching_database_database_2001}{}{teachingManagement/logical/teaching_database_1/generated_sql}{01_teaching_database_database_2001.sql}{}{}{}%
\listingSQLandOutput{2-}{teaching_1:01_teaching_database_database_2001}{}{0.05}{0.4}{0.9}{0.6}%
\end{frame}%
%
\begin{frame}[t]%
\frametitle{Tabelle student}%
\parbox{0.435\linewidth}{\small%
\begin{itemize}%
%
\item Hier sehen wir das auto-generierte Skript, um die Tabelle \sqlil{student} zu erstellen.%
%
\item<2-> Sie ist dafür da, die Entitäten vom Typ~\emph{Student} zu speichern.%
%
\item<3-> Sie hat den Ersatz-Primärschlüssel~\sqlil{id} und speichert die Namen der Studenten als variabel-langen String.%
%
\end{itemize}%
}%
\gitLoadAndExecSQL{teaching_1:03_public_student_table_5071}{}{teachingManagement/logical/teaching_database_1/generated_sql}{03_public_student_table_5071.sql}{teaching_database}{}{}%
\listingSQLandOutput{}{teaching_1:03_public_student_table_5071}{}{0.47}{0.2}{0.529}{0.87}
\end{frame}%
%
\begin{frame}[t]%
\frametitle{Tabelle professor}%
\parbox{0.435\linewidth}{\small%
\begin{itemize}%
%
\item Hier sehen wir das auto-generierte Skript, um die Tabelle \sqlil{professor} zu erstellen.%
%
\item<2-> Sie ist dafür da, die Entitäten vom Typ~\emph{Professor} zu speichern.%
%
\item<3-> Sie hat genau die gleiche Struktur, wie die Tabelle \sqlil{student}.%
%
\end{itemize}%
}%
\gitLoadAndExecSQL{teaching_1:04_public_professor_table_5087}{}{teachingManagement/logical/teaching_database_1/generated_sql}{04_public_professor_table_5087.sql}{teaching_database}{}{}%
\listingSQLandOutput{}{teaching_1:04_public_professor_table_5087}{}{0.47}{0.2}{0.529}{0.87}
\end{frame}%
%
\begin{frame}[t]%
\frametitle{Tabelle student}%
\parbox{0.435\linewidth}{\small%
\begin{itemize}%
%
\item Hier sehen wir das auto-generierte Skript, um die Tabelle \sqlil{module} zu erstellen.%
%
\item<2-> Sie ist dafür da, die Entitäten vom Typ~\emph{Module} zu speichern.%
%
\item<3-> Sie hat ebenfalls einen Ersatz-Primärschlüssel~\sqlil{id} und speichert die Titel der Module als variabel-langen String.%
%
\end{itemize}%
}%
\gitLoadAndExecSQL{teaching_1:05_public_module_table_5095}{}{teachingManagement/logical/teaching_database_1/generated_sql}{05_public_module_table_5095.sql}{teaching_database}{}{}%
\listingSQLandOutput{}{teaching_1:05_public_module_table_5095}{}{0.47}{0.2}{0.529}{0.87}
\end{frame}%
%
\begin{frame}[t]%
\frametitle{Tabelle course}%
\parbox{0.435\linewidth}{\small%
\begin{itemize}%
%
\only<-4>{%
\item Hier sehen wir das auto-generierte Skript, um die Tabelle \sqlil{course} zu erstellen.%
}%
%
\item<2-> Diese Tabelle ist etwas interessanter.%
%
\item<3-> Sie hat natürlich wieder einen Ersatz-Primärschlüssel~\sqlil{id}.%
%
\item<4-> Davon abgesehen speichert sie zwei Fremdschlüssel -- \sqlil{professor} und \sqlil{module} -- die auf die Tabellen mit den selben Namen zeigen~(was ein paar Slides weiter mit \sqlil{REFERENCES}-Einschränkungen erzwungen wird).%
%
\item<5-> Dazu speichern wir noch das Semester als einfache Ganzzahl.%
%
\item<6-> Wir werden die \inQuotes{Jahreszahl * 10} verwenden, gefolgt von einer 1 oder 2, je nachdem, ob wir im Frühlings- oder Herbstsemester sind.%
%
\end{itemize}%
}%
\gitLoadAndExecSQL{teaching_1:06_public_course_table_5105}{}{teachingManagement/logical/teaching_database_1/generated_sql}{06_public_course_table_5105.sql}{teaching_database}{}{}%
\listingSQLandOutput{}{teaching_1:06_public_course_table_5105}{}{0.47}{0.2}{0.529}{0.87}
\end{frame}%
%
\begin{frame}[t]%
\frametitle{Tabelle enrolls}%
\parbox{0.435\linewidth}{\small%
\begin{itemize}%
%
\only<-5>{%
\item Hier sehen wir das auto-generierte Skript, um die Tabelle \sqlil{enrolls} zu erstellen.%
}%
%
\item<2-> Diese Tabelle ist anders: Sie hat \alert{keinen} Ersatz-Primärschlüssel!%
%
\item<3-> Sie referenziert die Tabellen \sqlil{student} und \sqlil{course} über entsprechende Fremdschlüssel~(die wir nachher mit Einschränkungen durchsetzen).%
%
\item<4-> Diese beiden Fremdschlüssel bilden gleichzeitig den Primärschlüssel der Tabelle.%
%
\item<5-> Dadurch sind die Paare ihrer Wert automatisch \sqlil{UNIQUE}.%
%
\item<6-> Dadurch kann sich kein Student mehr als einmal in einen Kurs einschreiben.%
%
\end{itemize}%
}%
\gitLoadAndExecSQL{teaching_1:07_public_enrolls_table_5131}{}{teachingManagement/logical/teaching_database_1/generated_sql}{07_public_enrolls_table_5131.sql}{teaching_database}{}{}%
\listingSQLandOutput{}{teaching_1:07_public_enrolls_table_5131}{}{0.47}{0.2}{0.529}{0.87}
\end{frame}%
%
\begin{frame}[t]%
\frametitle{Einschränkungen~(1)}%
\parbox{0.435\linewidth}{\small%
\begin{itemize}%
%
\item Hier sehen wir das auto-generierte Skript, das die Einschränkung erstellt, die erzwingt, dass jede Zeile in Tabelle \sqlil{course} mit genau einer Zeile in Tabelle \sqlil{professor} verbunden ist.%
%
\end{itemize}%
}%
\gitLoadAndExecSQL{teaching_1:08_public_course_course_professor_fk_constraint_5116}{}{teachingManagement/logical/teaching_database_1/generated_sql}{08_public_course_course_professor_fk_constraint_5116.sql}{teaching_database}{}{}%
\listingSQLandOutput{}{teaching_1:08_public_course_course_professor_fk_constraint_5116}{}{0.47}{0.2}{0.529}{0.87}
\end{frame}%
%
\begin{frame}[t]%
\frametitle{Einschränkungen~(2)}%
\parbox{0.435\linewidth}{\small%
\begin{itemize}%
%
\item Hier sehen wir das auto-generierte Skript, das die Einschränkung erstellt, die erzwingt, dass jede Zeile in Tabelle \sqlil{course} mit genau einer Zeile in Tabelle \sqlil{module} verbunden ist.%
%
\end{itemize}%
}%
\gitLoadAndExecSQL{teaching_1:09_public_course_course_module_fk_constraint_5117}{}{teachingManagement/logical/teaching_database_1/generated_sql}{09_public_course_course_module_fk_constraint_5117.sql}{teaching_database}{}{}%
\listingSQLandOutput{}{teaching_1:09_public_course_course_module_fk_constraint_5117}{}{0.47}{0.2}{0.529}{0.87}
\end{frame}%
%
\begin{frame}[t]%
\frametitle{Einschränkungen~(3)}%
\parbox{0.435\linewidth}{\small%
\begin{itemize}%
%
\item Hier sehen wir das auto-generierte Skript, das die Einschränkung erstellt, die erzwingt, dass jede Zeile in Tabelle \sqlil{enrolls} mit genau einer Zeile in Tabelle \sqlil{student} verbunden ist.%
%
\end{itemize}%
}%
\gitLoadAndExecSQL{teaching_1:10_public_enrolls_enrolls_student_fk_constraint_5145}{}{teachingManagement/logical/teaching_database_1/generated_sql}{10_public_enrolls_enrolls_student_fk_constraint_5145.sql}{teaching_database}{}{}%
\listingSQLandOutput{}{teaching_1:10_public_enrolls_enrolls_student_fk_constraint_5145}{}{0.47}{0.2}{0.529}{0.87}
\end{frame}%
%
\begin{frame}[t]%
\frametitle{Einschränkungen~(4)}%
\parbox{0.435\linewidth}{\small%
\begin{itemize}%
%
\item Hier sehen wir das auto-generierte Skript, das die Einschränkung erstellt, die erzwingt, dass jede Zeile in Tabelle \sqlil{enrolls} mit genau einer Zeile in Tabelle \sqlil{course} verbunden ist.%
%
\end{itemize}%
}%
\gitLoadAndExecSQL{teaching_1:11_public_enrolls_enrolls_course_fk_constraint_5151}{}{teachingManagement/logical/teaching_database_1/generated_sql}{11_public_enrolls_enrolls_course_fk_constraint_5151.sql}{teaching_database}{}{}%
\listingSQLandOutput{}{teaching_1:11_public_enrolls_enrolls_course_fk_constraint_5151}{}{0.47}{0.2}{0.529}{0.87}
\end{frame}%
%
\begin{frame}[t]%
\frametitle{Daten Einfügen}%
\parbox{0.435\linewidth}{\small%
\begin{itemize}%
%
\only<-5>{%
\item Nun schreiben wir ein eigenes \glsShort{SQL}-Skript, um einige Daten in die Tabellen einzupflegen.%
}%
%
\only<-6>{%
\item<2-> Wir definieren die drei Studenten  Bibbo, Bebbo und Bebba.%
}%
%
\only<-7>{%
\item<3-> Als Professoren nehmen wir Weise und Bobbo.%
}%
%
\only<-8>{%
\item<4-> Wir erstellen drei Module: \emph{\pgls{python}}, \emph{Databases} und \emph{\pgls{Java}}.%
}%
%
\only<-9>{%
\item<5-> Über die Tabelle \sqlil{courses} spezifizieren wir, dass Prof.~Weise teaches \emph{\pgls{python}} und \emph{Databases} im Semester~20252 unterrichtet, also im Herbstsemester~2025.%
}%
%
\only<-11>{%
\item<6-> Prof.~Bobbo unterrichtet \emph{\pgls{Java}} im Frühlings- und Herbstsemester~2026.%
}%
%
\item<7-> Herr Bibbo schreibt sich in den \emph{\pgls{Java}}-Kurs von Prof.~Bobbo ein und in beide Kurse von by Prof.~Weise.%
%
\item<8-> Herr Bebbo nimmt stattdessen \emph{\pgls{Java}} und \emph{\pgls{python}}.%
%
\item<9-> Frau Bebba schreibt sich in \emph{\pgls{Java}} und \emph{Databases} ein.%
%
\item<10-> Mit 4~\sqlilIdx{INNER JOIN}s können wir alle Daten wieder zusammensetzen.%
%
\item<11-> Wir sortieren die Ergebnisse nach Studentenname, Professorname, Modultitel und Semester.%
%
\end{itemize}%
}%
\gitLoadAndExecSQL{teaching_database_1:insert_and_select}{}{teachingManagement/logical/teaching_database_1}{insert_and_select.sql}{teaching_database}{}{}%
\listingSQLandOutput{}{teaching_database_1:insert_and_select}{}{0.47}{0.01}{0.529}{0.99}
\end{frame}%
%
%
\begin{frame}[t]%
\frametitle{Aufräumen}%
\begin{itemize}%
\item Räumen wir nun auf.
\end{itemize}%
%
\gitLoadAndExecSQL{teaching_database_1:cleanup}{}{teachingManagement/logical/teaching_database_1}{cleanup.sql}{}{}{}%
\listingSQLandOutput{}{teaching_database_1:cleanup}{}{0.05}{0.3}{0.9}{0.7}%
\end{frame}%
%%
\section{Zusammenfassung}%
%
\begin{frame}%
\frametitle{Zusammenfassung: Beziehungen Höherer Ordnung}%
\begin{itemize}%
\item Wir haben gesehen, dass wir eine konzeptuelle ternäre Beziehung auf verschiedene Arten im logischen Schema implementieren können.%
%
\item<2-> Diesmal waren es zwei verschiedene Arten, aber es hätten genauso gut auch mehr sein können.%
%
\item<3-> Wir können uns hier aber nicht alle möglichen ternären Beziehungen anschauen.%
%
\item<4-> Die Frage, wie wir ternäre Beziehungen oder Beziehungen noch höherer Ordnung implementieren ist also nicht einfach oder erschöpfend zu beantworten.%
%
\item<5-> Man braucht einfach Erfahrung, um dafür ein Gefühl entwickeln zu können.%
%
\item<6-> So oder so wird man die Beziehungen oft auf mehrere binäre Beziehungen herunterbrechen.%
%
\item<7-> Und wie man die implementiert wissen wir ja.%
\end{itemize}%
\end{frame}%
%
\begin{frame}%
\frametitle{Zusammenfassung}%
\begin{itemize}%
\item An dieser Stelle haben wir eigentlich alles abgearbeitet, was man wissen muss, um konzeptuelle Schemas in logische Modelle zu transformieren.%
%
\item<2-> Wir können Entitäten zu Tabellen machen.%
%
\item<3-> Wir können Beziehungen implementieren.%
%
\item<4-> Wir kommen mit allen möglichen Randfällen wie Beziehungsattributen, schwachen Entitäten, und Beziehungen höherer Ordnung klar.%
%
\item<6-> Was uns noch fehlt sind ein paar grundlegende Richtlinien, wie man \emph{gute} logische Modelle erstellt.%
\end{itemize}%
\end{frame}%
%%
%%
\endPresentation%
\end{document}%%
\endinput%
%
